\documentclass[11pt]{article}

\usepackage[hmargin=0.95in,vmargin=0.85in]{geometry}
\usepackage[T1]{fontenc}
\usepackage{lmodern}
\usepackage{microtype}
\usepackage{amsmath,amssymb,amsthm,mathtools}
\usepackage{booktabs,array}
\usepackage{enumitem}
\usepackage{graphicx}
\usepackage{xcolor}
\usepackage[hidelinks]{hyperref}

\setlength{\parindent}{0pt}
\setlength{\parskip}{5pt}
\setlength{\emergencystretch}{3em}
\setlist[enumerate]{leftmargin=1.6em,itemsep=2pt,topsep=3pt}
\setlist[itemize]{leftmargin=1.45em,itemsep=2pt,topsep=3pt}

\newtheorem{proposition}{Proposition}
\newtheorem{lemma}{Lemma}
\newtheorem{corollary}{Corollary}
\theoremstyle{definition}
\newtheorem{definition}{Definition}
\newtheorem{assumption}{Assumption}
\theoremstyle{remark}
\newtheorem{remark}{Remark}

\newcommand{\dd}{\mathrm d}
\newcommand{\ind}{\mathbf 1}
\newcommand{\E}{\mathbb E}
\newcommand{\R}{\mathbb R}
\newcommand{\mcN}{\mathcal N}
\newcommand{\mcH}{\mathcal H}
\newcommand{\mcS}{\mathcal S}

\begin{document}

\begin{center}
{\LARGE\bfseries Housing Reallocation and Demographic Transitions}\par
\vspace{3pt}
{\Large Technical appendix for a two-generation OLG model}\par
\end{center}

\section{Environment and timing}

Time is discrete, and one period is one generation.  At date $t$, the economy
contains $Y_t$ young households and $O_t$ old households.  Each young household
has a permanent type $i=(y_i,b_i)$ drawn from a probability distribution $F$,
where $y_i$ is current income and $b_i$ is liquid wealth available at closing.
Each old household exits at the end of the period.  A child born at $t$
generates $\nu>0$ young households at $t+1$ after survival and household
formation.  Children do not inherit their parent's individual financial state.

A young household consumes $c$, occupies $h$ units of housing, and chooses
completed fertility $n$.  Its utility is
\begin{equation}
 u_t^Y(c,h,n)=\log c+\alpha\log(h-\kappa n)+\vartheta_t\log n.
 \label{eq:young_u_app}
\end{equation}
Each child uses $\chi>0$ units of current goods and $\kappa>0$ units of space.
An old household consumes $c^O$, retains $h^O$ units of housing, and leaves an
estate $e$ valued according to
\begin{equation}
 u^O(c^O,h^O,e)=\log c^O+\gamma\log h^O+\omega_B\log e.
 \label{eq:old_u_app}
\end{equation}
All choices lie in the log domain.  In particular, $c,n,h-\kappa n,c^O,h^O$,
and $e$ are strictly positive.

The economy has a fixed housing stock $\bar H$.  Each unit produces one unit of
housing services in every period.  Renters can occupy at most $\bar h^R$, while
owners can occupy at most $\bar h^O>\bar h^R$.  The house price is $P_t$.  A
claim to one unit of date-$t+1$ goods costs $q\in(0,1)$ at date $t$, and the
corresponding gross return is $R=q^{-1}$.

A competitive housing intermediary buys a unit at $t$, pays the recurring
property tax $\tau^pP_t$, rents the unit for $r_t$, and sells it at $t+1$.
Zero profit therefore requires
\begin{equation}
 r_t=(1+\tau^p)P_t-qP_{t+1}.
 \label{eq:rent_app}
\end{equation}
A young buyer finances the share $\phi\in(0,1)$ of the purchase price with a
one-period mortgage, posts $(1-\phi)P_th$ at closing, and repays $R\phi P_th$
when old.  A household sale at $t$ realizes the per-unit tax
\begin{equation}
 \ell_t=\tau^g[P_t-P_{t-1}]_+.
 \label{eq:gains_app}
\end{equation}
Housing retained until death receives a stepped-up basis and is liquidated by
the estate without this household capital-gains liability.  Intermediaries
mark their basis to market.  Old owners self-occupy retained housing and do not
operate as landlords.

The cash flows imply two private user costs.  A young buyer who sells the
marginal unit when old pays
\begin{equation}
 r_t^O=(1+\tau^p)P_t-q(P_{t+1}-\ell_{t+1})
      =r_t+q\ell_{t+1}.
 \label{eq:buyer_cost_app}
\end{equation}
An old incumbent who retains a marginal unit instead of selling it pays
\begin{equation}
 r_t^I=(P_t-\ell_t)+\tau^pP_t-qP_{t+1}
      =r_t-\ell_t.
 \label{eq:incumbent_cost_app}
\end{equation}
The buyer anticipates the tax on next-period resale, while the incumbent avoids
the current realization tax by retaining the unit.

The timing within date $t$ is: inherited states and $P_{t-1}$ are observed;
$P_t$, $P_{t+1}$, $T_t$, and $T_{t+1}$ are correctly anticipated; young
households choose tenure, housing, fertility, consumption, and saving; old
households choose sales, retained housing, consumption, and estates; markets
clear and taxes are rebated; the young become the next old and their children
form the next young cohort.

\section{Household problems and exact conditional solutions}

\subsection{Old renters and old incumbents}

An old renter with financial wealth $a$ solves
\begin{align}
 V_t^R(a)=\max_{c^O,h^O,e}\;&
 \log c^O+\gamma\log h^O+\omega_B\log e
 \label{eq:old_renter_app}\\
 \text{subject to}\quad&
 c^O+q e+r_th^O=a+T_t,\qquad 0<h^O\leq\bar h^R.\nonumber
\end{align}
An old incumbent enters with net financial wealth $a$ and inherited housing
$H$.  Selling all housing produces $(P_t-\ell_t)H$.  Expressing the estate in
total goods, including retained housing, gives
\begin{align}
 V_t^O(a,H)=\max_{c^O,h^O,e}\;&
 \log c^O+\gamma\log h^O+\omega_B\log e
 \label{eq:old_owner_app}\\
 \text{subject to}\quad&
 c^O+q e+(r_t-\ell_t)h^O
 =a+(P_t-\ell_t)H+T_t,\nonumber\\
 &0<h^O\leq H,\qquad e\geq P_{t+1}h^O.\nonumber
\end{align}
The last inequality is the nonnegative liquid-estate restriction.  It prevents
the household from promising an estate smaller than the value of the retained
house.

Let
\begin{equation}
 A=1+\gamma+\omega_B,\qquad
 \rho_t=r_t-\ell_t,\qquad
 X_t=a+(P_t-\ell_t)H+T_t.
 \label{eq:A_rho_X}
\end{equation}
The old problems are strictly concave on their feasible sets.  Their active
sets therefore provide a complete characterization.

\begin{lemma}[Old-age allocation]
Suppose $r_t>0$, $\rho_t>0$, and resources are sufficient for the log domain.
If the renter size limit is slack, then
\begin{equation}
 c^O=\frac{a+T_t}{A},\qquad
 h^O=\frac{\gamma(a+T_t)}{Ar_t},\qquad
 e=\frac{\omega_B(a+T_t)}{Aq}.
 \label{eq:old_renter_closed}
\end{equation}
If both incumbent constraints are slack, then
\begin{equation}
 c^O=\frac{X_t}{A},\qquad
 h^O=\frac{\gamma X_t}{A\rho_t},\qquad
 e=\frac{\omega_BX_t}{Aq}.
 \label{eq:old_owner_closed}
\end{equation}
The two slackness conditions are
\begin{equation}
 \frac{\gamma X_t}{A\rho_t}<H,
 \qquad
 \omega_B\rho_t>\gamma qP_{t+1}.
 \label{eq:old_slackness}
\end{equation}
\end{lemma}

\begin{proof}
For a generic budget $c^O+qe+p_hh^O=X$, the three first-order conditions imply
$p_hh^O=\gamma c^O$ and $qe=\omega_Bc^O$.  The budget then gives
$Ac^O=X$.  Substituting $p_h=r_t$ for a renter and $p_h=\rho_t$ for an
incumbent gives equations~\eqref{eq:old_renter_closed} and
\eqref{eq:old_owner_closed}.  Substituting the incumbent solution into
$h^O<H$ and $e>P_{t+1}h^O$ gives equation~\eqref{eq:old_slackness}.
\end{proof}

For completeness, the remaining incumbent active sets also have closed forms.
If only estate composition binds, define $d_t=\rho_t+qP_{t+1}$; then
\begin{equation}
 c^O=\frac{X_t}{A},\qquad
 h^O=\frac{(\gamma+\omega_B)X_t}{Ad_t},\qquad
 e=P_{t+1}h^O.
 \label{eq:old_estate_binding}
\end{equation}
If only retention binds, let $Z_t=X_t-\rho_tH$; then
\begin{equation}
 h^O=H,\qquad
 c^O=\frac{Z_t}{1+\omega_B},\qquad
 e=\frac{\omega_BZ_t}{(1+\omega_B)q},
 \label{eq:old_retention_binding}
\end{equation}
provided the estate inequality remains slack.  If both constraints bind,
$h^O=H$, $e=P_{t+1}H$, and
\begin{equation}
 c^O=X_t-(\rho_t+qP_{t+1})H.
 \label{eq:old_both_binding}
\end{equation}
The old renter has the analogous size-cap solution: set $h^O=\bar h^R$ and
allocate the remaining resources between $c^O$ and $e$ in shares
$1:(\omega_B)$.
The appropriate row is selected by feasibility and complementary slackness.
For example, the estate-binding row requires
$\gamma qP_{t+1}\geq\omega_B\rho_t$; the retention-only row requires
\begin{equation}
 \omega_BZ_t>(1+\omega_B)qP_{t+1}H,\qquad
 \gamma Z_t\geq(1+\omega_B)\rho_tH.
 \label{eq:retention_only_conditions}
\end{equation}
When both constraints bind, the two multipliers are nonnegative exactly when
\begin{equation}
 \frac{q}{c^O}\geq\frac{\omega_B}{P_{t+1}H},
 \qquad
 \frac{\gamma+\omega_B}{H}\geq\frac{\rho_t+qP_{t+1}}{c^O}.
 \label{eq:both_old_multipliers}
\end{equation}

In the fully interior incumbent region, the indirect value and its wealth
envelope are
\begin{align}
 V_t^O(a,H)
 &=A\log X_t-A\log A+\gamma\log\gamma-\gamma\log\rho_t
   +\omega_B\log\omega_B-\omega_B\log q,
 \label{eq:old_indirect}\\
 V_{a,t}^O&=\frac{A}{X_t}=\frac{1}{c_t^O},\qquad
 V_{H,t}^O=(P_t-\ell_t)V_{a,t}^O.
 \label{eq:old_envelopes}
\end{align}
These envelopes are the link between the young buyer's current problem and
her old-age resources.

\subsection{Young renters and young owners}

Conditional on renting, a young household solves
\begin{align}
 W_t^R(i)=\max_{c,h,n,a'}\;&u_t^Y(c,h,n)+\beta V_{t+1}^R(Ra')
 \label{eq:young_renter_app}\\
 \text{subject to}\quad&
 c+\chi n+a'+r_th=y_i+b_i+T_t,\nonumber\\
 &0<h\leq\bar h^R,\qquad a'\geq0.\nonumber
\end{align}
Conditional on ownership, it solves
\begin{align}
 W_t^O(i)=\max_{c,h,n,a'}\;&u_t^Y(c,h,n)
 +\beta V_{t+1}^O(a_{t+1},h)
 \label{eq:young_owner_app}\\
 \text{subject to}\quad&
 c+\chi n+a'+[(1-\phi)+\tau^p]P_th=y_i+b_i+T_t,\nonumber\\
 &a_{t+1}=Ra'-R\phi P_th,\nonumber\\
 &(1-\phi)P_th\leq b_i,\qquad
 0<h\leq\bar h^O,\qquad a'\geq0.\nonumber
\end{align}
Liquid wealth $b_i$, rather than income received after closing, is pledgeable
for the down payment.

The first-order conditions also characterize cases in which saving or a
housing limit binds.  Let $\lambda_t^m$ be the multiplier on the current budget,
$\sigma_t^m$ the multiplier on $a'\geq0$, $\eta_t^m$ the multiplier on the
tenure-specific size limit, and $\mu_t$ the owner down-payment multiplier.  The
common consumption and fertility conditions are
\begin{align}
 \lambda_t^m&=\frac1{c_t^m},
 \label{eq:young_lambda}\\
 \frac{\vartheta_t}{n_t^m}
 &=\lambda_t^m\chi+
   \frac{\alpha\kappa}{h_t^m-\kappa n_t^m}.
 \label{eq:young_n_kkt}
\end{align}
The saving conditions are
\begin{align}
 \lambda_t^R&=\beta R V_{a,t+1}^R+\sigma_t^R,
 \label{eq:renter_saving_kkt}\\
 \lambda_t^O&=\beta R V_{a,t+1}^O+\sigma_t^O,
 \label{eq:owner_saving_kkt}
\end{align}
with $\sigma_t^ma_t'=0$.  When saving and the future old constraints are slack,
the housing conditions reduce to
\begin{equation}
 \frac{\alpha c_t^m}{h_t^m-\kappa n_t^m}
 =r_t^m+\zeta_t^m,
 \qquad r_t^R=r_t,\quad r_t^O=r_t+q\ell_{t+1},
 \label{eq:young_housing_kkt}
\end{equation}
where the current-goods access values are
\begin{equation}
 \zeta_t^R=\frac{\eta_t^R}{\lambda_t^R},\qquad
 \zeta_t^O=
 \underbrace{\frac{\mu_t}{\lambda_t^O}(1-\phi)P_t}_{\zeta_t^{O,F}}
 +\underbrace{\frac{\eta_t^O}{\lambda_t^O}}_{\zeta_t^{O,S}}.
 \label{eq:zeta_app}
\end{equation}
Thus $\zeta_t^{O,F}$ is the financing component and $\zeta_t^{O,S}$ is the
owner-size component.  If both limits bind, choices identify their sum but not
the two multipliers separately.

\subsection{Lifetime reduction and unconstrained policies}

Suppose young saving and the future old allocation are interior.  The saving
Euler equation and equation~\eqref{eq:old_owner_closed} imply
\begin{equation}
 \frac1{c_t}=\beta R\frac1{c_{t+1}^O},\qquad
 c_{t+1}^O=\beta Rc_t,\qquad
 qX_{t+1}=kc_t,\qquad k\equiv\beta A.
 \label{eq:euler_reduction}
\end{equation}
For an owner,
\begin{equation}
 X_{t+1}=Ra_t'-R\phi P_th_t+(P_{t+1}-\ell_{t+1})h_t+T_{t+1}.
 \label{eq:owner_old_resources}
\end{equation}
Combining current and discounted old resources gives the lifetime identity
\begin{equation}
 (1+k)c_t+\chi n_t+r_t^mh_t=\mathcal W_{it},\qquad
 \mathcal W_{it}=y_i+b_i+T_t+qT_{t+1}.
 \label{eq:lifetime_app}
\end{equation}
The renter identity is the same with $r_t^m=r_t$; the owner identity uses
$r_t^m=r_t+q\ell_{t+1}$.  This is a private lifetime budget, not an aggregate
resource constraint.

\begin{proposition}[Unconstrained conditional policies]
Let
\begin{equation}
 B_t=1+k+\alpha+\vartheta_t.
 \label{eq:B_app}
\end{equation}
When saving, housing limits, and future old constraints are all slack, the
conditional young policies under tenure $m$ are
\begin{align}
 c_t^m&=\frac{\mathcal W_{it}}{B_t},
 \label{eq:c_uncon}\\
 n_t^m&=\frac{\vartheta_t\mathcal W_{it}}
 {B_t(\chi+\kappa r_t^m)},
 \label{eq:n_uncon}\\
 h_t^m-\kappa n_t^m&=\frac{\alpha\mathcal W_{it}}{B_tr_t^m},
 \label{eq:s_uncon}\\
 h_t^m&=\frac{\alpha\mathcal W_{it}}{B_tr_t^m}
 +\frac{\kappa\vartheta_t\mathcal W_{it}}
 {B_t(\chi+\kappa r_t^m)}.
 \label{eq:h_uncon}
\end{align}
The associated saving policies are
\begin{align}
 a_t^{\prime R}&=kc_t^R-qT_{t+1},
 \label{eq:saving_r_closed}\\
 a_t^{\prime O}&=kc_t^O-qT_{t+1}
 -[q(P_{t+1}-\ell_{t+1})-\phi P_t]h_t^O.
 \label{eq:saving_o_closed}
\end{align}
\end{proposition}

\begin{proof}
Let $s=h-\kappa n$ denote family space.  The unconstrained housing condition
gives $s=\alpha c/r_t^m$.  Combining the fertility condition with this result
gives $n=\vartheta_t c/(\chi+\kappa r_t^m)$.  Substituting
$h=s+\kappa n$ in equation~\eqref{eq:lifetime_app} gives
$\mathcal W_{it}=B_tc$.  Equations~\eqref{eq:saving_r_closed} and
\eqref{eq:saving_o_closed} follow from $qX_{t+1}=kc_t$.
\end{proof}

The effective housing limits are
\begin{equation}
 \widehat h_t^R=\bar h^R,\qquad
 \widehat h_{it}^O=\min\left\{\bar h^O,
 \frac{b_i}{(1-\phi)P_t}\right\}.
 \label{eq:effective_caps}
\end{equation}
A limit binds exactly when the unconstrained policy in
equation~\eqref{eq:h_uncon} exceeds the relevant $\widehat h$.  When it binds,
$h=\widehat h$ and fertility is the unique root of
\begin{equation}
 F(n,\widehat h)equiv
 \frac{\vartheta_t}{n}
 -\frac{\alpha\kappa}{\widehat h-\kappa n}
 -\frac{(1+k)\chi}
 {\mathcal W_{it}-r_t^m\widehat h-\chi n}=0.
 \label{eq:binding_cap_root}
\end{equation}
The root lies in
\begin{equation}
 0<n<\min\left\{\frac{\widehat h}{\kappa},
 \frac{\mathcal W_{it}-r_t^m\widehat h}{\chi}\right\}.
 \label{eq:binding_cap_domain}
\end{equation}
Since $F_n<0$, the root is unique.  Consumption and saving then follow from
equations~\eqref{eq:lifetime_app}, \eqref{eq:saving_r_closed}, and
\eqref{eq:saving_o_closed}.  If $a'=0$, the lifetime reduction no longer
applies; equations~\eqref{eq:young_lambda}--\eqref{eq:owner_saving_kkt} and the
binding resource constraint uniquely characterize the conditional solution.

The household chooses $m_t(i)\in\{R,O\}$ by comparing the two conditional
values.  Conditional problems are strictly concave, but the tenure decision is
discrete.  The aggregate analysis below uses the following standard
regularity condition.

\begin{assumption}[Tenure regularity]
The type distribution is atomless at tenure indifference: for every equilibrium
price and transfer vector,
$F\{i:W_t^R(i)=W_t^O(i)\}=0$.  Conditional policies are unique, and parameters
remain in a compact region in which resources and relevant user costs are
strictly positive.
\label{ass:tenure_regular}
\end{assumption}

Under Assumption~\ref{ass:tenure_regular}, individual tenure switches do not
create jumps in aggregate choice maps.  A continuous idiosyncratic tenure taste
is one primitive way to satisfy the condition, but it is not needed for the
household derivations.

\section{Aggregation and positive steady states}

\subsection{States, distributions, and government accounts}

Let $G_t$ be the normalized distribution of old financial wealth, inherited
housing, and tenure.  It is generated by the preceding young policies:
\begin{equation}
 G_{t+1}=\operatorname{Law}_{i\sim F}\!\left(
 Ra_t'(i)-\ind\{m_t(i)=O\}R\phi P_th_t(i),
 \ind\{m_t(i)=O\}h_t(i),m_t(i)\right).
 \label{eq:G_pushforward_app}
\end{equation}
For every bounded measurable test function $\psi$, this statement means
\begin{align}
 \int\psi(a,H,m)\,\dd G_{t+1}
 =\int \psi\big(&Ra_t'(i)-\ind\{m_t(i)=O\}R\phi P_th_t(i),
 \nonumber\\[-4pt]
 &\ind\{m_t(i)=O\}h_t(i),m_t(i)\big)\,\dd F(i).
 \label{eq:G_test_app}
\end{align}
Because children do not inherit individual states, this is a one-generation
pushforward rather than an intergenerational wealth fixed point.

Let $\bar n_t$, $\bar h_t^Y$, and $\bar h_t^O$ be average fertility, young
housing, and old housing.  Let
\begin{equation}
 \bar s_t^O=\int \ind\{m=O\}\,[H-h_t^O(a,H,m)]\,\dd G_t
 \label{eq:old_sales_app}
\end{equation}
be average taxable old-owner sales, assigning zero to renters.  Cohorts and
housing satisfy
\begin{equation}
 Y_{t+1}=\nu\bar n_tY_t,\qquad O_{t+1}=Y_t,\qquad
 Y_t\bar h_t^Y+O_t\bar h_t^O=\bar H.
 \label{eq:aggregate_laws_app}
\end{equation}
The government rebates property- and gains-tax revenue equally across living
households:
\begin{equation}
 (Y_t+O_t)T_t=\tau^pP_t\bar H+\ell_tO_t\bar s_t^O.
 \label{eq:gov_app}
\end{equation}
Property taxation applies to the fixed occupied stock.  The gains-tax base is
the housing actually sold by old owners.

\subsection{A two-equation characterization of a positive steady state}

Fix a constant fertility preference $\vartheta$.  Let
$\mcN(P,T;\vartheta)$ denote average young fertility and let
$\mcH^Y(P,T;\vartheta)$ and $\mcH^O(P,T;\vartheta)$ denote stationary housing
per young and old household after solving the household problems and applying
the pushforward in equation~\eqref{eq:G_pushforward_app}.  Write
\begin{equation}
 \mcH(P,T;\vartheta)
 =\mcH^Y(P,T;\vartheta)+\mcH^O(P,T;\vartheta).
 \label{eq:Hsum_app}
\end{equation}

\begin{definition}[Positive steady state]
A positive steady state is a constant equilibrium with $Y^*=O^*>0$, constant
$P^*$ and $T^*$, stationary policies and $G^*$, and all household, housing,
cohort, and government conditions satisfied.
\end{definition}

In any steady state, $P_t=P_{t-1}=P_{t+1}=P^*$, and therefore
\begin{equation}
 \ell^*=0,\qquad r^*=\delta P^*,\qquad
 \delta\equiv1+\tau^p-q>0.
 \label{eq:ss_user_cost_app}
\end{equation}
The absence of real capital gains at a stationary endpoint follows from the tax
base itself; it does not require a separate assumption.

\begin{proposition}[Reduced steady-state system]
Every positive steady state is characterized by a root $(P^*,T^*)\in\R_{++}^2$
of
\begin{align}
 F_1(P,T;\vartheta)
 &\equiv\nu\mcN(P,T;\vartheta)-1=0,
 \label{eq:ss_F1}\\
 F_2(P,T;\vartheta)
 &\equiv T-\frac{\tau^pP}{2}\mcH(P,T;\vartheta)=0.
 \label{eq:ss_F2}
\end{align}
Given the root, the population scale and old distribution are
\begin{equation}
 Y^*=O^*=\frac{\bar H}{\mcH(P^*,T^*;\vartheta)},
 \qquad G^*=\Gamma(P^*,T^*;\vartheta).
 \label{eq:ss_scale_G}
\end{equation}
Consequently, every positive steady state has average fertility
$\mcN(P^*,T^*;\vartheta)=1/\nu$.
\end{proposition}

\begin{proof}
Stationarity of $O_{t+1}=Y_t$ gives $O^*=Y^*$.  Stationarity of the young cohort
then gives $1=\nu\mcN$.  Housing clearing gives
$\bar H=Y^*\mcH$.  Since $\ell^*=0$, the government budget gives
$2Y^*T^*=\tau^pP^*\bar H$.  Eliminating $Y^*$ using housing clearing produces
equation~\eqref{eq:ss_F2}; reversing these steps establishes sufficiency.
\end{proof}

The positive steady state is therefore a property of primitives and policy,
not a normalization.  It need not exist for arbitrary parameters.  The next
result provides transparent sufficient conditions.

\begin{proposition}[Existence and local uniqueness]
Suppose Assumption~\ref{ass:tenure_regular} holds on the rectangle
$\mathcal D=[\underline P,\overline P]\times
[\underline T,\overline T]\subset\R_{++}^2$.  Suppose further that:
\begin{enumerate}
 \item the fiscal map
 $\Psi(P,T)=(\tau^pP/2)\mcH(P,T;\vartheta)$ maps
 $[\underline T,\overline T]$ into itself for every $P$ and is a uniform
 contraction in $T$;
 \item its unique fixed point $T=\mathcal T(P)$ is interior and continuous;
 and
 \item the replacement residual changes sign:
 \begin{equation}
 F_1(\underline P,\mathcal T(\underline P);\vartheta)
 F_1(\overline P,\mathcal T(\overline P);\vartheta)<0.
 \label{eq:ss_sign_change}
 \end{equation}
\end{enumerate}
Then a positive steady state exists in $\mathcal D$.  It is locally unique if
the Jacobian
\begin{equation}
 J=\begin{bmatrix}
 \nu\mcN_P & \nu\mcN_T\\[2pt]
 -\dfrac{\tau^p}{2}(\mcH+P\mcH_P) &
 1-\dfrac{\tau^pP}{2}\mcH_T
 \end{bmatrix}_{(P^*,T^*;\vartheta)}
 \label{eq:ss_J}
\end{equation}
is nonsingular.
\end{proposition}

\begin{proof}
The contraction theorem gives a unique continuous fiscal locus
$T=\mathcal T(P)$.  The intermediate value theorem applied to
$F_1(P,\mathcal T(P);\vartheta)$ gives existence.  Nonsingularity of $J$ and
the implicit-function theorem give local uniqueness.
\end{proof}

When the Jacobian is nonsingular, a change in the fertility preference has the
local comparative static
\begin{equation}
 \begin{bmatrix}P_\vartheta^*\\ T_\vartheta^*\end{bmatrix}
 =-J^{-1}
 \begin{bmatrix}
 \nu\mcN_\vartheta\\[2pt]
 -\dfrac{\tau^pP^*}{2}\mcH_\vartheta
 \end{bmatrix}.
 \label{eq:ss_ift}
\end{equation}
The population-scale response then follows from
\begin{equation}
 \frac{\dd\log Y^*}{\dd\vartheta}
 =-\frac{
 \mcH_P P_\vartheta^*+\mcH_T T_\vartheta^*+\mcH_\vartheta
 }{\mcH}.
 \label{eq:ss_scale_derivative}
\end{equation}
Equations~\eqref{eq:ss_ift}--\eqref{eq:ss_scale_derivative} show why replacement
fertility can be invariant across endpoints while prices, transfers, individual
policies, the old-state distribution, and population scale differ.

\section{Perfect-foresight transition equilibrium}

Suppose $\vartheta_t=\vartheta^0$ for $t<0$ and
$\vartheta_t=\vartheta^1$ for $t\geq0$, and that both values admit positive
steady states $\mathcal E^0$ and $\mathcal E^1$.  The predetermined aggregate
state is
\begin{equation}
 S_t=(Y_t,O_t,G_t,P_{t-1}).
 \label{eq:predetermined_state}
\end{equation}
The lagged price records the basis of the old owners who may sell at $t$.
Prices and transfers dated $t$ are not predetermined.

\begin{definition}[Transition equilibrium]
Given $S_0=(Y^0,Y^0,G^0,P^0)$, a perfect-foresight transition equilibrium is a
sequence $\{P_t,T_t,S_t\}_{t\geq0}$ and household policies such that:
\begin{enumerate}
 \item every old and young household solves its problem under the correctly
 anticipated sequence;
 \item $G_{t+1}$ is the pushforward in equation~\eqref{eq:G_pushforward_app};
 \item housing clears and the government budget balances at every date;
 \item cohorts evolve according to equation~\eqref{eq:aggregate_laws_app};
 \item $\ell_t$ and $r_t$ satisfy equations~\eqref{eq:gains_app} and
 \eqref{eq:rent_app}; and
 \item $(P_t,T_t,S_t)$ converges to $\mathcal E^1$.
\end{enumerate}
\end{definition}

This definition initializes one inherited state, not two shocks.  The date-zero
old are those generated by the old steady state, and their basis is $P_{-1}=P^0$.
The permanent preference change is the only exogenous innovation.

To display the equilibrium operator, let household solutions under a candidate
perfect-foresight sequence generate
$\bar h_t^Y$, $\bar h_t^O$, $\bar n_t$, and $\bar s_t^O$.  At each date the two
equilibrium residuals are
\begin{align}
 \mathcal R_t^H
 &=Y_t\bar h_t^Y+O_t\bar h_t^O-\bar H,
 \label{eq:transition_housing_resid}\\
 \mathcal R_t^G
 &=(Y_t+O_t)T_t-\tau^pP_t\bar H-\ell_tO_t\bar s_t^O.
 \label{eq:transition_gov_resid}
\end{align}
Together with the forward laws for $S_{t+1}$, a transition is a root of the
sequence operator $\mathcal R=\{\mathcal R_t^H,\mathcal R_t^G\}_{t\geq0}$.

The capital-gains rule can equivalently be written as the complementarity
system
\begin{equation}
 \ell_t\geq0,\qquad
 \ell_t\geq\tau^g(P_t-P_{t-1}),\qquad
 \ell_t\,[\ell_t-\tau^g(P_t-P_{t-1})]=0.
 \label{eq:gains_complementarity}
\end{equation}
The steady state lies at the kink.  Ordinary derivatives are therefore
regime-specific: the appreciating regime has derivative $\tau^g$, while the
non-appreciating regime has derivative zero.

\begin{proposition}[Finite-horizon existence and local determinacy]
Fix a terminal date $J$ and impose the terminal continuation
$\mathcal E^1$ from $J+1$ onward.  Suppose aggregate policies are continuous,
the induced state law remains in a compact invariant set, and the mapping from
a candidate vector
$(P_0,\ldots,P_J,T_0,\ldots,T_J)$ to the vector that clears
equations~\eqref{eq:transition_housing_resid}--\eqref{eq:transition_gov_resid}
is a continuous self-map of a compact convex set.  Then a finite-horizon
transition exists.  Within a fixed tenure, constraint, and appreciation regime,
it is locally unique if the Jacobian of the stacked residual vector with
respect to the price-transfer vector is nonsingular.
\end{proposition}

\begin{proof}
Brouwer's theorem gives a fixed point of the continuous clearing map.  Within a
fixed active set, the residual vector is differentiable.  A nonsingular stacked
Jacobian and the implicit-function theorem then give local uniqueness.
\end{proof}

The proposition states sufficient, verifiable conditions; it does not assert
that every parameterization has a transition.  At an appreciation kink one
must either verify the candidate active regime, solve the complementarity
problem directly, or use a generalized-Jacobian argument.  For an infinite
horizon, a sequence of finite-horizon solutions converges to a transition if
it is uniformly bounded and has a convergent subsequence whose limit satisfies
the terminal condition and equilibrium equations.

\begin{corollary}[One shock can produce gains for many cohorts]
Following the single permanent preference shock, if the equilibrium price
satisfies $P_s>P_{s-1}$ for $s=t,\ldots,t+J$, then
$\ell_s>0$ on all $J+1$ dates.  A two-period household life does not restrict
capital gains to one date.
\end{corollary}

\begin{proof}
Equation~\eqref{eq:gains_app} gives
$\ell_s=\tau^g(P_s-P_{s-1})>0$ on each appreciating date.  Two-period lives
determine whose basis is $P_{s-1}$; they do not determine the number of dates on
which the equilibrium price can appreciate.
\end{proof}

\section{Intergenerational housing allocation}

\subsection{Marginal values}

Consider a young owner $i$ whose down-payment constraint binds, whose owner-size
limit is slack, and whose saving and next-period old constraints are slack.
Consider an old incumbent $j$ whose retention and estate-composition constraints
are slack and for whom $r_t>\ell_t$.  Their marginal values of one unit of
current housing services, in date-$t$ goods, are
\begin{equation}
 MV_{it}^Y=r_t+q\ell_{t+1}+\zeta_{it}^{O,F},
 \qquad
 MV_{jt}^O=r_t-\ell_t.
 \label{eq:pair_values_app}
\end{equation}

\begin{lemma}[Transition valuation gap]
Under the stated active sets,
\begin{equation}
 MV_{it}^Y-MV_{jt}^O
 =\zeta_{it}^{O,F}+\ell_t+q\ell_{t+1}.
 \label{eq:valuation_gap_app}
\end{equation}
\end{lemma}

\begin{proof}
Subtract the incumbent user cost in equation~\eqref{eq:incumbent_cost_app} from
the buyer's constrained marginal value in
equation~\eqref{eq:young_housing_kkt}.  The common rental user cost cancels.
\end{proof}

The three terms have distinct sources.  The financing wedge is the value of
pledgeable liquidity at closing.  The current gains tax lowers the incumbent's
private cost of retention.  The discounted future gains tax raises the buyer's
cost of acquiring a unit that will be sold when old.

\subsection{Steady-state compensated reallocation}

The welfare comparison first holds fertility, cohort masses, prices, tenure,
and all other household choices fixed.  It changes only a marginal housing unit
and balanced transfers between the two living households.

\begin{definition}[Local compensated housing test]
An allocation passes the local compensated housing test if no infinitesimal
transfer of title and occupancy from an old incumbent to a current young owner,
together with feasible balanced goods transfers, can leave both weakly better
off and one strictly better off.  The current young household's lifetime utility
includes her own old age.
\end{definition}

Let $K_{ij}\geq0$ be the real date-$t$ resource cost per unit of executing and
servicing the reallocation.

\begin{proposition}[Steady-state reallocation]
At a positive steady state, the goods-equivalent surplus from transferring a
marginal unit from old incumbent $j$ to financing-constrained young owner $i$
is
\begin{equation}
 \mcS_{ij}^*=\zeta_i^{O,F}-K_{ij}.
 \label{eq:ss_surplus_app}
\end{equation}
If $\mcS_{ij}^*>0$, the steady-state allocation fails the local compensated
housing test.
\end{proposition}

\begin{proof}
At a steady state, $\ell_t=\ell_{t+1}=0$, so the valuation gap is
$\zeta_i^{O,F}$.  Transfer $\varepsilon>0$ units of title and occupancy from
$j$ to $i$.  To preserve the buyer's binding closing constraint, lend
\begin{equation}
 L=(1-\phi)P^*\varepsilon
 \label{eq:bridge_ss}
\end{equation}
immediately before closing and require repayment $RL$ when the buyer is old.
The bridge has zero present value.  Let $\lambda_i$ and $\mu_i$ be the buyer's
budget and down-payment multipliers.  Its contribution to the buyer's
Lagrangian is
\begin{equation}
 (\lambda_i+\mu_i)L-\beta V_{a,i,t+1}^ORL
 =\mu_iL=\lambda_i\zeta_i^{O,F}\varepsilon,
 \label{eq:bridge_value_ss}
\end{equation}
where $\lambda_i=\beta RV_{a,i,t+1}^O$ is the saving Euler equation.  Thus the
bridge relaxes the closing constraint without creating net financial wealth.

The incumbent sells the same marginal unit.  At an interior retention choice,
net sale proceeds and the induced estate adjustment compensate the first-order
loss of housing services because $u_h^O/u_c^O=r^*$.  The fixed property-tax base
is unchanged, and no gains tax is realized at a constant price.  After paying
$K_{ij}\varepsilon$, the coalition has
$(\zeta_i^{O,F}-K_{ij})\varepsilon$ in date-$t$ goods.  If this amount is
positive, a sufficiently small transfer can be divided so that neither
household loses and one gains strictly.
\end{proof}

The proposition is a sufficient local result.  A nonpositive surplus for one
pair does not establish global efficiency: another pair, active set, or
reallocation instrument may still generate gains.

\subsection{A two-date transition result}

Along a transition, the tax terms in equation~\eqref{eq:valuation_gap_app}
change government rebates.  A Pareto statement therefore requires an explicit
two-date transfer technology rather than only a comparison of private marginal
values.

\begin{assumption}[Two-date compensation technology]
A reallocation authority can execute a taxable marginal sale, issue the
zero-present-value pre-closing bridge in equation~\eqref{eq:bridge_ss}, trade
one-period claims at price $q$, require the transferred marginal title to be
resold at $t+1$, make balanced lump-sum side transfers at $t$ and $t+1$,
reclaim only incremental tax rebates, and pay real cost $K_{ijt}$ in date-$t$
goods.
\label{ass:two_date_transfer}
\end{assumption}

\begin{proposition}[Transition reallocation]
Under Assumption~\ref{ass:two_date_transfer} and the active sets used in
equation~\eqref{eq:pair_values_app}, the present-value surplus from a marginal
transfer is
\begin{equation}
 \mcS_{ijt}=\zeta_{it}^{O,F}+\ell_t+q\ell_{t+1}-K_{ijt}.
 \label{eq:transition_surplus_app}
\end{equation}
If $\mcS_{ijt}>0$, the transfer supports a local two-date constrained Pareto
improvement.
\end{proposition}

\begin{proof}
Transfer $\varepsilon$ units.  The bridge loan produces the financing gain
$\zeta_{it}^{O,F}\varepsilon$ by the same Euler-equation cancellation as in
equation~\eqref{eq:bridge_value_ss}.  The old incumbent is compensated to first
order by the sale proceeds and estate adjustment implied by her interior
retention condition.

The extra sale realizes $\ell_t\varepsilon$ in date-$t$ tax revenue.  The
required next-period resale is feasible without changing the buyer's active set
because her retention constraint is slack; it realizes
$\ell_{t+1}\varepsilon$ at $t+1$.  If $N_s=Y_s+O_s$, the equal-rebate rule
implies
\begin{equation}
 N_t\dd T_t=\ell_t\varepsilon,
 \qquad
 N_{t+1}\dd T_{t+1}=\ell_{t+1}\varepsilon.
 \label{eq:incremental_rebates}
\end{equation}
The authority reclaims exactly each household's incremental rebate, so every
other household remains at its baseline allocation.  This yields
$\ell_t\varepsilon+q\ell_{t+1}\varepsilon$ in date-$t$ present value.  Charging
the buyer no more than its bridge gain, paying $K_{ijt}\varepsilon$, and
distributing the strictly positive residual leaves every affected household
weakly better off and at least one strictly better off.  Differentiability and
a strict first-order inequality imply a genuine improvement for sufficiently
small $\varepsilon$.
\end{proof}

The compensation experiment fixes fertility and population.  If a policy
changes the number or identity of future households, ordinary Pareto comparisons
do not by themselves rank the two populations.  One must state an
endogenous-population welfare criterion, such as the $P$- or $A$-efficiency
concepts developed by \cite{golosovjonestertilt2007}; see also
\cite{pereznievasetal2019}.  No such criterion is needed for Propositions 5 and 6
because every household whose resources change is held weakly above its
baseline and cohort masses are fixed.

\section{Housing access and fertility}

The reallocation result concerns willingness to pay for a fixed housing unit.
The fertility result instead asks how an access limit changes a household's
optimal $n$ while tenure, prices, lifetime wealth, and the active set are fixed.

Suppose $h=\widehat h$ binds and fertility remains interior.  Define
\begin{equation}
 D=\mathcal W_{it}-r_t^m\widehat h-\chi n,\qquad
 S=\widehat h-\kappa n,\qquad c=\frac{D}{1+k}.
 \label{eq:D_S_app}
\end{equation}
Implicit differentiation of equation~\eqref{eq:binding_cap_root} gives the exact
partial-equilibrium derivative
\begin{equation}
 \frac{\dd n}{\dd\widehat h}
 =\frac{
 \dfrac{\alpha\kappa}{S^2}
 -\dfrac{(1+k)\chi r_t^m}{D^2}}
 {\dfrac{\vartheta_t}{n^2}
 +\dfrac{(1+k)\chi^2}{D^2}
 +\dfrac{\alpha\kappa^2}{S^2}}.
 \label{eq:fertility_derivative_app}
\end{equation}
The denominator is strictly positive.  More family space raises the return to
fertility, represented by the first numerator term; the goods cost of occupying
additional housing lowers disposable consumption, represented by the second.

At a constrained allocation, equation~\eqref{eq:young_housing_kkt} implies
\begin{equation}
 \frac{D}{S}=\frac{(1+k)(r_t^m+\zeta_t^m)}{\alpha}.
 \label{eq:D_S_ratio}
\end{equation}
For $r_t^m>0$, substitution gives the exact sign condition
\begin{equation}
 \frac{\dd n}{\dd\widehat h}\geq0
 \quad\Longleftrightarrow\quad
 \chi\leq
 \frac{(1+k)\kappa(r_t^m+\zeta_t^m)^2}{\alpha r_t^m}.
 \label{eq:fertility_sign_app}
\end{equation}
A stronger condition stated only in primitives and the user cost is
\begin{equation}
 \chi\leq\frac{(1+k)\kappa r_t^m}{\alpha}.
 \label{eq:fertility_sufficient_app}
\end{equation}
The sign is therefore economic, not automatic: a child must be sufficiently
space intensive relative to its goods cost.

When the down-payment limit is the active owner constraint,
\begin{equation}
 \widehat h_{it}^O=\frac{b_i}{(1-\phi)P_t},
 \label{eq:finance_cap_app}
\end{equation}
and its mechanical derivatives are
\begin{equation}
 \frac{\partial\widehat h^O}{\partial b_i}
 =\frac1{(1-\phi)P_t},\qquad
 \frac{\partial\widehat h^O}{\partial\phi}
 =\frac{b_i}{(1-\phi)^2P_t},\qquad
 \frac{\partial\widehat h^O}{\partial P_t}
 =-\frac{\widehat h^O}{P_t}.
 \label{eq:cap_pass_through}
\end{equation}
For a reform that changes only the pledgeable cap locally,
$\partial n/\partial z=(\dd n/\dd\widehat h)
(\partial\widehat h/\partial z)$.  A funded policy or a general-equilibrium
change also changes $\mathcal W$, $r_t^m$, rebates, and possibly tenure; those
additional channels must be included before assigning an aggregate sign.

\section{A solved equilibrium path}

This section gives a reproducible equilibrium example of the preceding
equations.  It uses the owner-only restriction of the model, a binding
down-payment limit, interior saving, and interior old choices.  The example is
a constructive existence check for this restricted parameterization rather
than an empirical calibration.

\begin{table}[ht]
\centering
\caption{Parameters of the worked equilibrium}
\label{tab:worked_parameters}
\small
\begin{tabular}{@{}llcllc@{}}
\toprule
Object & Parameter & Value & Object & Parameter & Value\\
\midrule
Discount price & $q$ & 0.50 & Mortgage share & $\phi$ & 0.80\\
Property tax & $\tau^p$ & 0.05 & Gains-tax rate & $\tau^g$ & 0.20\\
Space weight & $\alpha$ & 0.40 & Space per child & $\kappa$ & 0.50\\
Goods per child & $\chi$ & 0.15 & Young discount & $\beta$ & 0.40\\
Old housing weight & $\gamma$ & 0.30 & Estate weight & $\omega_B$ & 0.40\\
Income & $y$ & 1.00 & Liquid wealth & $b$ & 0.06\\
Household formation & $\nu$ & 2.00 & Housing stock & $\bar H$ & 2.00\\
\bottomrule
\end{tabular}
\end{table}

At date zero, $\vartheta$ rises permanently from $0.35$ to $0.50$.  Solving
equations~\eqref{eq:ss_F1}--\eqref{eq:ss_F2} gives the two endpoints in
Table~\ref{tab:worked_steady_states}.

\begin{table}[ht]
\centering
\caption{Steady-state endpoints of the worked equilibrium}
\label{tab:worked_steady_states}
\small
\begin{tabular}{@{}lrr@{}}
\toprule
Object & Initial steady state & Final steady state\\
\midrule
Fertility preference $\vartheta$ & 0.3500 & 0.5000\\
Average fertility $\bar n$ & 0.5000 & 0.5000\\
House price $P$ & 0.3999 & 0.5599\\
Per-household rebate $T$ & 0.0130 & 0.0130\\
Young housing $h^Y$ & 0.7503 & 0.5358\\
Old housing $h^O$ & 0.5453 & 0.3894\\
Young and old cohort mass $Y=O$ & 1.5438 & 2.1617\\
\bottomrule
\end{tabular}
\end{table}

Average fertility equals $1/\nu=0.5$ at both endpoints, as required by
replacement.  The endpoints nevertheless differ in prices, individual housing,
and population scale.  On impact, fertility rises above replacement.  The
young cohort consequently expands, and the house price appreciates while the
fixed stock is reallocated across changing cohort masses.

\begin{figure}[ht]
\centering
\includegraphics[width=0.80\textwidth]{figures/simplified_olg_transition_example.png}
\caption{Perfect-foresight response to one permanent increase in the fertility
preference.  Dotted and dashed lines in the price panel denote the initial and
final steady-state prices.  The gains tax is positive on each of the first four
dates because the equilibrium price continues to appreciate.  Fertility moves
on impact, cohort masses respond with the OLG lag, and both return to a positive
steady state.}
\label{fig:worked_transition}
\end{figure}

The finite-horizon root uses the final steady state as the continuation value
after date 24.  The maximum absolute housing or government residual is
$3.2\times10^{-10}$, and the maximum terminal-state gap is
$2.4\times10^{-10}$.  Along the path, the minimum old sale is $0.145$, the
minimum incumbent user cost is $0.206$, the minimum estate-composition slack is
$0.029$, minimum saving is $0.408$, and the minimum difference between
unconstrained owner housing and the down-payment cap is $0.348$.  The active
sets used in the analytical derivation therefore hold strictly in the computed
path.  Replacing the terminal horizon 24 with horizon 20 changes the first nine
dates by at most $1.4\times10^{-15}$, and the analytical fertility derivative
in equation~\eqref{eq:fertility_derivative_app} matches a centered
finite-difference derivative to within $2.0\times10^{-10}$.

The example establishes the narrow transition point mechanically.  A single
permanent preference change moves fertility on impact, cohort sizes with a lag,
and prices over several dates.  Capital gains are a transition object and
vanish at both steady-state endpoints.

\section{Result map}

The model delivers four separate conclusions.
\begin{enumerate}
 \item A positive steady state, when it exists, is a root of the two equations
 in $(P,T)$ given by equations~\eqref{eq:ss_F1}--\eqref{eq:ss_F2}.  Replacement
 fixes average fertility but not prices, policies, distributions, or population
 scale.
 \item A single permanent shock can generate a multi-date price transition and
 capital gains for several successive old cohorts.  A second shock is not
 required.
 \item At a steady state, a financing wedge supports a local compensated
 housing reallocation when $\zeta_i^{O,F}>K_{ij}$.  Along a transition, the
 analogous two-date surplus additionally contains $\ell_t+q\ell_{t+1}$, provided
 the rebate and transfer ledger in Assumption~\ref{ass:two_date_transfer} is
 available.
 \item Relaxing a housing-access limit raises fertility under the exact
 condition in equation~\eqref{eq:fertility_sign_app}.  This response is distinct
 from the Pareto comparison and is not, by itself, a general-equilibrium welfare
 result.
\end{enumerate}

\begin{thebibliography}{9}
\small
\setlength{\itemsep}{2pt}

\bibitem[Golosov, Jones, and Tertilt(2007)]{golosovjonestertilt2007}
Golosov, Mikhail, Larry E. Jones, and Mich\`ele Tertilt. 2007.
``Efficiency with Endogenous Population Growth.''
\emph{Econometrica} 75(4): 1039--1071.

\bibitem[P\'erez-Nievas et al.(2019)]{pereznievasetal2019}
P\'erez-Nievas, Mikel, Jos\'e I. Conde-Ruiz, and Eduardo L. Gim\'enez. 2019.
``Efficiency and Endogenous Fertility.''
\emph{Theoretical Economics} 14(2): 475--512.

\end{thebibliography}

\end{document}
